Apply the simple-graph five-cycle counts to the red edges of
$G_{\rm s}$. Completing the simple support to a $D$-regular simple
graph adds only blue edges. All vertices of a red five-cycle and of
its five parts are incident with red edges. They therefore belong
to the old graph, which is induced in the completion. Thus the
completion preserves both the red five-cycles and their five
parts. Write $q_C$ for the number of ordered induced red five-cycles
divided by $2sD^3$ and $L_R$ for the sum of $\frac RD$ over these
cycles divided by the same quantity. If each such cycle satisfies
$R\geq\alpha D$, then summing this inequality gives $L_R\geq\alpha q_C$.

Applying the simple-graph certificate inequality to four choices
of objective gives
\begin{equation}\label{eq:support-projection-four}
\begin{aligned}
 S-0.99999q_C
   &\le \frac{5}{4}-\frac{0.99999}{8}+2\cdot10^{-9},\\
 S+0.5682L_R
   &\le \frac{5}{4}+10^{-9},\\
 S+0.4897(L_R-0.32q_C)
   &\le \frac{5}{4}-\frac{0.4897\cdot0.32}{8}+2\cdot10^{-9},\\
 S+0.45119\left(L_R-\frac{q_C}{2}\right)
   &\le \frac{5}{4}-\frac{0.45119}{16}+2\cdot10^{-9}.
\end{aligned}
\end{equation}
The coefficients and the error allowances in
\eqref{eq:support-projection-four} are rational.
\textbf{Appendix~\ref{app:multi-projection}} specifies the
four certificate checks. To see why the displayed coefficients
suffice, write each certified inequality as $S+\gamma Y\leq
\frac54+\varepsilon$, moving the multiple of $\frac18$ to the left.
For $Y\geq0$, replacing $\gamma$ by a smaller nonnegative coefficient
preserves the inequality. For $Y<0$ the bound follows from
$S\leq\frac54$.

For a given $\alpha>0$, substitute $L_R\geq\alpha q_C$ in
\eqref{eq:support-projection-four}. Each of the four inequalities
now has the form $S\leq a_i-b_iq_C$. Include $S\leq\frac54$ as
a fifth inequality, with $a_5=\frac54$ and $b_5=0$. Thus
\[
 S\leq S_{\max}(\alpha):=
 \max_{x\geq0}\min_{1\leq i\leq5}(a_i-b_ix),
 \qquad e\geq1-\frac{S_{\max}(\alpha)}{M}.
\]
The second line of \eqref{eq:support-projection-four} has
$b_2=0.5682\alpha>0$. The minimum of the five lines is continuous
and piecewise linear and tends to $-\infty$ as $x$ tends to infinity.
Its maximum is therefore attained at zero or at the nonnegative
intersection of two lines with different slopes. This gives a finite
rational calculation for $S_{\max}(\alpha)$.

These inequalities can also be imposed directly on $X$. Multiplying
each simple-graph inequality by $1-e=\frac{s}{m}$ changes the cycle
normalization from $2sD^3$ to $2mD^3$. Its term $(1-e)S$ is $O_0$,
and the same factor changes $L_R$ to the corresponding weighted
coordinate sum. All terms of the resulting inequality are linear
in the weighted counts.

Let $k\geq3$ and let $K$ be the vertex set of a red $K_k$ in the weighted host.
For each vertex $v$ put $a_v=|N_{G_{\rm s}}(v)\cap K|$ and let
$h_v$ and $d_{\rm blue}(v)$ count the marked and the unmarked edges
incident with $v$, respectively. Let $m_K$ count the marked edges with
both endpoints in $K$, with multiplicity, and put
$J_K=\sum_{v\in K}j_v$. For a marked edge $uv$, the set
$(N_{G_{\rm s}}(u)\cup N_{G_{\rm s}}(v))\cap K$ meets every edge
of the red clique. Indeed, a disjoint red edge in $K$ must have an
endpoint adjacent to $u$ or $v$, and the same holds when the red
edge shares an endpoint with $uv$. Every vertex cover of $K_k$ has
size at least $k-1$, so $a_u+a_v\geq k-1$. If $u\in K$, then
$a_u=k-1$ and $a_v\geq1$. If both endpoints lie in $K$, then
$a_u+a_v=2(k-1)$. Summing over all marked edges with multiplicity gives
\[
 (k-1)m+\sum_{v\in K}h_v+(k-3)m_K
 \le\sum_v a_vh_v
 \le D(kD-J_K)-\sum_{v\in K}d_{\rm blue}(v).
\]
For the upper bound, $\sum_v a_v=\sum_{v\in K}(D-j_v)=kD-J_K$
in the regular completion. Moreover, $a_vd_{\rm blue}(v)$ is at
least the number of blue neighbors of $v$ in $K$. Summing this
last inequality counts each blue incidence at $K$ and gives the
upper bound on $\sum_v a_vh_v$, since $h_v=D-d_{\rm blue}(v)$.
Using $\sum_{v\in K}(h_v+d_{\rm blue}(v))=kD$ now gives
\[
 (k-1)m\le kD(D-1)-DJ_K-(k-3)m_K.
\]
For $k=5$, the right side is at most $5D(D-1)-20$. Thus $4m\leq5D(D-1)-20$ if a red $K_5$ exists.
This is impossible when $m\geq\frac54D^2$. For odd $D$ it is also
impossible when $m=f(D)+1$, since then $4m=5D^2-2D+5$.
The coordinates of every pattern containing a red $K_5$ are
therefore zero in either case.

Fix a red edge $xy$. Let $U=N_{G_{\rm s}}(x)\setminus (N_{G_{\rm s}}(y)\cup\{y\})$
and let $W$ be the vertices adjacent to neither $x$ nor $y$. In the
host induced on $U\cup W$, mark only the previously marked edges
between $U$ and $W$. Edges within a part can still join endpoints of
these marked edges. For each marked edge $uv$ between $U$ and $W$, and
each red neighbor $z$ of $y$, the strong clique condition for $uv$
and $yz$ forces $uz$ or $vz$ to be a host edge. All such vertices $z$ lie
outside $U\cup W$, and different $z$ give different edges incident
with $u$ or $v$ that are absent from the induced host. Hence the Ore-degree $\sigma$ of the bipartite
marked subgraph in this induced host is at most $2D$ minus the
number of red neighbors of $y$.

If there is a marked edge between $U$ and $W$, the corresponding
upper bounds are
\[
 \frac{2973}{10000}\sigma^2\quad\text{for all copies},\qquad
 \frac{273323}{10^6}\sigma^2\quad\text{for base copies},\qquad
 \frac14\sigma^2\quad\text{for extra copies}.
\]
We give the counts that yield these bounds. In the induced host
let $H$ be the bipartite graph of marked edges between $U$ and $W$,
with their multiplicities. Suppose $H$ is nonempty, put $h=\Delta(H)$, and take
a vertex $v$ of marked degree $h$. Let $A=N_H(v)$,
$B'=V\setminus N[v]$ (where $V$ and $N[v]$ refer to the induced
host), $p=|A|$ and $\ell=\sigma-d(v)$, with $d(v)$ the degree in
the host. Every $a\in A$ has host degree at most $\ell$ since $va$
is marked. Let $T$ be the number of marked edges meeting $A$,
$P$ the number of unmarked edges from $A$ to $B'$, and
$F=H[B']$. Write $f=|E(F)|$ and $u=\Delta(F)$, with $u=0$ if
$F$ is edgeless. Every marked edge omitted by the counts $T$ and
$f$ has an endpoint in $N(v)\setminus A$. This set has at most
$d(v)-h$ vertices. Consequently
\[
 |E(H)|\leq h(\sigma-h)-h\ell+T+f,\qquad T+P\leq p\ell .
\]
For each $a\in A$, its neighbors in $B'$ meet every edge of $F$.
Since neither endpoint of that edge is adjacent to $v$, one must
be adjacent to $a$ to join the edge to $va$. Thus $f\leq u\ell$.
The same observation shows that the Ore-degree of the marked edges
in the host induced on $B'$ is at most $\sigma-p$. Indeed, for
each of the $p$ vertices of $A$, at least one edge to an endpoint
of the marked edge is lost when the host is restricted to $B'$.

First count all marked copies. On each pair joined by a marked
edge, remove the unmarked copies. On every other adjacent pair
retain one unmarked edge. This preserves all marked edges and all
adjacencies, and $\sigma$ still bounds the degree sum on each marked
edge. Use the counts above for this reduced host. Let $X'$ comprise the vertices of
$B'$ in the same part of $H$ as $v$. For $x\in X'$, let $h_x$
count the marked edges from $x$ to $A$, with multiplicity, and put
$f_x=d_F(x)$.
Each marked edge of $F$ has one endpoint in $X'$. For every
$a\in A$ it has an endpoint adjacent to $a$. The adjacency to $a$
either has an unmarked copy or is marked and meets the endpoint
in $X'$. Counting these possibilities yields
\[
 pf\leq uP+\sum_{x\in X'} h_xf_x .
\]
If $T,f>0$, let $\alpha=\frac{T}{T+f}$ and $\gamma=\frac{f}{T+f}$.
The identity
$(h_x+f_x)(\gamma^2h_x+\alpha^2 f_x)-h_x f_x
=(\gamma h_x-\alpha f_x)^2$ and the bound $h_x+f_x\leq h$
give $\sum_{x\in X'} h_xf_x\leq h(\gamma^2T+\alpha^2 f)
=\frac{hTf}{T+f}$. Substituting the resulting lower bound for
$P$ in $T+P\leq p\ell$ gives
\[
 |E(H)|\leq h(\sigma-h)+\left(1-\frac hp\right)T
       +\left(1-\frac hu\right)f
       +\frac{h^2Tf}{pu(T+f)} .
\]
If $f=0$, the previous counts give $|E(H)|\leq h(\sigma-h)$.
The induced host on $B'$ has fewer marked edges, since it omits
all marked edges at $v$. Applying bounds for this smaller graph with
Ore-degree at most $\sigma-p$ gives upper bounds for $f$.
The proof notes in the repository~\cite{FGSTCode} give further necessary
inequalities by comparing the marked neighborhoods in $A$ of
two vertices of $X'$. For example, if $x,x'\in X'$ are nonadjacent,
every vertex of $N_H(x)\cap A$ that is not adjacent to $x'$ must
be adjacent to every vertex of $N_H(x')\cap A$ that is not adjacent
to $x$. Otherwise two marked edges would have no edge joining their
endpoints. The rational interval check in \textbf{Appendix~\ref{app:ore-replay}}
uses these neighborhood inequalities and the bounds for the smaller
graph. It excludes $|E(H)|>\frac{2973}{10000}\sigma^2$ over the full
range of $\frac{h}{\sigma},\frac ph,\frac uh,\frac{T}{\sigma^2}$
and $\frac{f}{\sigma^2}$.

For the base copies, pass to the simple support, mark one edge on
each marked pair, and take the host induced on their endpoints.
This preserves every adjacency that joins two disjoint marked edges.
Let $s'$ be the number of marked edges in this simple host and let
$\sigma'\leq\sigma$ be their Ore-degree. Every vertex has a marked
neighbor, so its degree is at most $\sigma'-1$. For a connected
pattern $P$ on $k$ positions rooted at an ordered marked edge,
normalize its map count by $2s'(\sigma')^{k-2}$. The resulting
coordinate $q_P$ lies in $[0,1]$ by the spanning-tree bound and
vanishes if the marked edges of $P$ contain an odd cycle.
The count of two ordered marked edges in \eqref{eq:objective-value}
now gives the objective $\frac{s'}{(\sigma')^2}$.
The extension sum at one position is at most $q_P$. If $ij$ is a
marked edge of $P$, the two extension sums at $i$ and $j$ together
are at most $q_P$, since the sum of the two image degrees is at
most $\sigma'$. The moment matrices remain positive semidefinite
by \textbf{Lemma~\ref{lem:moment}}. The rational certificate in
\textbf{Appendix~\ref{app:ore-replay}} combines the extension inequalities
with nonnegative multipliers and the moment matrices with
positive semidefinite matrices. Each negative residual coefficient
is charged its absolute value, using $0\leq q_P\leq1$. The resulting
upper bound on $\frac{s'}{(\sigma')^2}$ is
$\frac{27332270490572378277}{10^{20}}<\frac{273323}{10^6}$.
This proves the asserted bound for base copies.

Finally, mark only the extra copies between $U$ and $W$, reserving
one base copy on each originally marked pair as an unmarked edge.
Every adjacent pair of vertices now has an unmarked edge joining it.
Apply the same counts with $H$ the graph of these extra copies.
The result is immediate if $H$ is empty. Otherwise, for
every pair consisting of $a\in A$ and an edge of $F$, the
joining adjacency has an unmarked copy, and each such copy serves
at most $u$ edges of $F$. Therefore $pf\leq uP$.
If $f=0$, the first two inequalities give $|E(H)|\leq
h(\sigma-h)$ because $p\leq h$. If $f>0$, then $u>0$, and
\[
 |E(H)|\leq h(\sigma-h)-(h-p)\ell+
          \left(1-\frac pu\right)f
 \leq h(\sigma-h)+\left(1-\frac hu\right)f
 \leq h(\sigma-h).
\]
Here the second inequality uses $f\leq u\ell$ and $p\leq h$,
and the last uses $u\leq h$. Finally $h(\sigma-h)\leq
\frac{\sigma^2}{4}$.

Multiplying any of the three bounds by the nonnegative forest
weight of the root and summing over root maps gives a linear
inequality for $X$. Expanding a neighbor count uses one further
position, and expanding its square or counting another edge uses
two.

Write $AX=(1,0,\ldots,0)^{\mathsf T}$ for the degree and sampling
equations, $NX\le0$ for the linear inequalities, and $G_i(X)\succeq0$
for the matrix inequalities. To bound $M$, take a vector
$\lambda$ with first entry $\lambda_0$, a nonnegative vector $\rho$, and matrices $P_i$
given as nonnegative sums of rational rank-one outer products.
Let $G_i^*$ be the adjoint defined by
$(G_i^*(P_i))^{\mathsf T}X=\tr(P_iG_i(X))$. The constraints give $\rho^{\mathsf T}NX\leq0$ and
$\tr(P_iG_i(X))\geq0$. Thus, for every feasible $X$,
\[
 \left(A^{\mathsf T}\lambda+N^{\mathsf T}\rho
        -\sum_iG_i^*(P_i)\right)^{\mathsf T}X
 \le\lambda_0 .
\]
Let $r$ be the coefficient vector on the left minus the coefficient
vector of $M$. If $0\le X_j\le U_j$ for each coordinate $j$, then
$M\le\lambda_0+\sum_{j:r_j<0}(-r_j)U_j$. For a coordinate $X(P,T)$,
an upper bound is
$\frac{1-e}{|\operatorname{Aut}(P,T)|}$ for $T=\varnothing$ and
$\frac{e}{|\operatorname{Aut}(P,T)|}$ otherwise. To obtain it, start at
a red edge, in $T$ if $T$ is nonempty, and extend along a
spanning tree containing $T$. Each new position contributes at
most $D$ in support count or extra multiplicity. The certificates
also use bounds on sums of coordinates of equal order and
equal forest size, obtained by counting all labeled spanning
trees containing the specified red root. Positive residuals
can be dropped since the coordinates are nonnegative. With these bounds the total correction is affine in $e$, so
its maximum over the allowed interval is attained at an endpoint. For an infeasibility certificate, the objective
is zero and the compensated right side is negative.
